\section{Introduction}
\label{sec:introduction}

Human action recognition in unconstrained broadcast video constitutes one of the foundational challenges in computer vision, with critical applications in automated athletic performance analytics, broadcasting highlight synthesis, biomechanical risk mitigation, and augmented refereeing~\cite{carreira2017i3d,feichtenhofer2019slowfast}. Among these domains, sports action recognition is uniquely demanding. Athletic movements are distinguished by rapid acceleration profiles, dynamic camera panning and zooming, intense body rotations, complex athletic apparatus interactions (e.g., rackets, vaults, gymnastics rings), and severe inter-class visual overlap~\cite{soomro2012ucf101}. For instance, distinguishing between \textit{BaseballPitch} and \textit{CricketBowling}, or separating \textit{TennisSwing} from \textit{TableTennisShot}, requires the computational model to decouple athletic kinematics from environmental clutter while capturing micro-second decisive moments such as ball release or racket impact.

Historically, video action recognition evolved through several architectural paradigms. Early methods coupled 2D convolutional networks (e.g., ResNet, ConvNeXt) with recurrent sequence models such as Long Short-Term Memory (LSTM) networks or Gated Recurrent Units (GRUs) to process temporally aggregated visual representations~\cite{donahue2015long}. While computationally tractable on commodity hardware, such decoupled models suffer from a fundamental bottleneck: spatial representations are computed frame-by-frame in isolation from temporal dynamics, preventing the joint optimization of spatiotemporal features. Subsequent advancements shifted toward 3D Convolutional Neural Networks (3D CNNs), including I3D~\cite{carreira2017i3d}, SlowFast~\cite{feichtenhofer2019slowfast}, and X3D~\cite{feichtenhofer2020x3d}. Although 3D CNNs achieve joint spatiotemporal convolution, their fixed receptive fields scale poorly across long temporal sequences, and their high parameter counts and cubic FLOP scaling create severe memory bottlenecks on budget-constrained hardware.

Recently, the advent of Vision Transformers (ViT)~\cite{dosovitskiy2020vit} and Video Masked Autoencoders (VideoMAE)~\cite{tong2022videomae,wang2023videomaev2} has established a new paradigm in visual representation learning. Unlike supervised backbones that risk overfitting to static background textures (e.g., green tennis courts or swimming pools), VideoMAE employs self-supervised pre-training with an aggressive tubelet masking ratio of 90\%. This extreme masking forces the network to infer missing spatiotemporal tokens by modeling deep motion vectors and dynamic body kinematics rather than memorizing static scene appearance. When fine-tuned, VideoMAE provides rich, generalizable representations of athletic movement.

Despite these algorithmic advances, real-world deployment and empirical validation of video foundation models face two major obstacles:
\begin{enumerate}
    \item \textbf{Data Leakage in Video Benchmarks}: In existing benchmarks such as UCF101~\cite{soomro2012ucf101}, video clips frequently originate from shared raw camera footage (e.g., $v\_[Class]\_g[Group]\_c[Clip].avi$). Naive random clip splitting inadvertently allocates clips from the same source group across both training and test partitions, resulting in severe visual scene contamination and falsely inflated test accuracies.
    \item \textbf{Hardware and Free-Tier Cloud Constraints}: Full fine-tuning of 86M-parameter vision transformers requires substantial GPU memory. Researchers relying on free-tier cloud accelerators (such as Kaggle and Google Colab) regularly face unpredictable session timeouts, disk space exhaustion, and out-of-memory (OOM) failures on 16~GB VRAM GPUs.
\end{enumerate}

To overcome these challenges, this paper presents an end-to-end Video Foundation Model framework for fine-grained sports action recognition across 82 categories. Our principal contributions are summarized as follows:
\begin{itemize}
    \item We formulate an end-to-end fine-tuning framework using **VideoMAE ViT-B** tailored for fine-grained sports action recognition across 82 classes (9,151 clips), utilizing $16 \times 16 \times 2$ spatiotemporal tubelet patch embeddings and layer-wise learning rate decay.
    \item We design and enforce a strict **Zero-Leakage Group-Wise Partitioning** scheme (70\% train, 15\% validation, 15\% test) based on source video group identifiers, guaranteeing completely disjoint scene, actor, and camera distributions to establish authentic generalization benchmarks.
    \item We develop a high-throughput, **interrupt-resilient cloud training engine** supporting PyTorch Distributed Data Parallelism (DDP) across dual-T4 GPUs on Kaggle and persistent Google Drive checkpointing on Colab, enabling zero-loss state recovery across session cuts.
    \item We benchmark VideoMAE against an extensive **4-tier baseline hierarchy** (Frozen Linear Probe, X3D-S, SlowFast-R50, and Video Swin-T) and provide comprehensive ablation studies across temporal frame density, spatial resolution, adaptation depth, and data scaling efficiency.
\end{itemize}
